% Ad Atticum 10.4 --- headnote
% ad-atticum-10-04, 15 April 49 BC, Cumae

\textit{Cicero to Atticus, written from the Cumaean villa on
the seventeenth day before the Kalends of May 49 BC (the
manuscript dateline: \textit{Scr.\ in Cumano xvii K. Mai.\
a.\ 705 (49)}). The letter opens as a reply to a packet of
Atticus's letters, one of them ``the size of a volume,''
and quickly turns to the two grievances that govern it: the
two ``supreme commanders'' --- Caesar and Pompey --- whom
Cicero now ranks beneath himself in fortune and in
conscience, on the doctrine he had laid down in his own
philosophical books that nothing is good but the honourable
and nothing evil but the base; and his nephew Quintus the
younger, whose journey to Caesar (after a meeting with
Hirtius) and whose denunciation of Cicero's plan to leave
Italy have just been reported to Cicero by what channel he
does not say. Section 5 carries the daggered crux
\textit{non tam} where the manuscripts plainly stumble. The
boy's father, Cicero's brother Quintus, ``lies broken in
grief''; the uncle, with characteristic legal precision,
asks Atticus to lay no blame on either kinsman, since the
trouble is the boy's nature --- the same nature, he says,
that ruined Curio and the younger Hortensius.}

\textit{Sections 7--11 are the letter's second movement: a
day-of dispatch from the visit Curio paid him at Cumae on
the Ides of April. Curio --- the tribune-turned-Caesarian
who had crossed the Rubicon with Caesar and was now en
route to take over Sicily --- speaks with his usual candour
and tells Cicero, in effect, everything: the Pompeian
exiles will be recalled; the Spanish provinces will fall to
Caesar; Caesar will then take an army to wherever Pompey
is; Pompey's death will be the end of \textit{[Greek: of
the matter]} (the manuscript daggers \textit{illi}, which
the translation marks); a great slaughter was narrowly
averted when Caesar lost his temper at Metellus the
tribune; Caesar himself is not cruel by nature but only as
clemency suits him politically. The vignette is one of the
most valuable inside views of the Caesarian camp that
survives. Curio also volunteers that Cicero may sail to
Greece openly through his province, which is the practical
gain Cicero records at the end of section 10. The closing
paragraph turns to domestic instructions: refer the Oppii
to Terentia; advise on the route to Rhegium; look after
Tiro. There are no Greek phrases in this letter, but three
daggered cruxes survive in the transmitted text (sections
5, 8, and 11), each marked with \textdagger\ in the
translation.}
